\chapter{Colocation Products and How They Are Sold}\label{ch:nw:colocation-products-and-how-they-are-sold}

European rules require a trading venue that sells co-location to give every user of the same service access to its network under the
same conditions, and the rule's list of conditions includes, between cooling and access to data, the length of the cable. In New
Jersey, Nasdaq has told the SEC that it is rebuilding the cabling of its data-centre campus so that every telecom carrier's connection
reaches its customers over an equal distance, and it has named the work the ``Equalization Project''. The customer whose cabinet
stands next to the matching engine and the customer at the far end of the hall are meant to wait the same nanoseconds for the same
bits, and where the room does not make the cables equal, a coil of fibre does.

Part~III of this book is about where the exchanges are and what it costs to be near them. This chapter describes what a venue's
data centre sells and on what terms: space, power, cooling and hands; \omterm{def:nw:colocation-products-and-how-they-are-sold:mmr}{cross-connects} and the \omterm{def:nw:colocation-products-and-how-they-are-sold:mmr}{meet-me room}; the fairness rules; the
connections into the venue's own network; and the fee schedules, which in the United States are public filings.

\section{What colocation sells: space, power, cooling and hands}

\begin{definition}[Colocation, proximity hosting]\label{def:nw:colocation-products-and-how-they-are-sold:colo}
\emph{Colocation}\index{colocation} is the rental of space, power and cooling for a customer's equipment in the data centre that houses a
trading venue's matching engines, with connections to the venue's network. \emph{Proximity hosting}\index{proximity hosting} is the same
service offered in a nearby data centre that the venue does not run, under an arrangement with the venue or with a third party.
\end{definition}

\begin{definition}[Colocation cabinet, colocation cage, power density]\label{def:nw:colocation-products-and-how-they-are-sold:cabinet}
A \emph{colocation cabinet}\index{colocation cabinet} is a lockable rack rented to one customer with a fixed allowance of power. A
\emph{colocation cage}\index{colocation cage} is a fenced area of the hall rented to one customer, holding several cabinets and access
controlled separately. The \emph{power density}\index{power density} of a cabinet is the power, in kilowatts, it is allowed to draw; it
sets its price as much as its space does.
\end{definition}

A venue's data centre is an industrial building. ICE's Mahwah data centre, which houses the NYSE exchanges, is a building of 398\,000
square feet on a 28-acre site 34 miles from Wall Street, with 28 megawatts of electrical power; its \omterm{def:nw:colocation-products-and-how-they-are-sold:cabinet}{colocation cabinets} come in 4, 8 and
12~kilowatts, scalable to 15, with partial cabinets of 1 to 2~kilowatts, each fed by two independent power circuits. Cabinets are sold by
\omterm{def:nw:colocation-products-and-how-they-are-sold:cabinet}{power density} because power and the cooling that removes it are the building's scarce resources: chapter~8's arithmetic of servers per
cabinet is the customer's side of the same constraint.

\begin{definition}[Remote hands]\label{def:nw:colocation-products-and-how-they-are-sold:hands}
\emph{Remote hands}\index{remote hands} is the data-centre operator's service of performing physical work on a customer's equipment at its
request (replacing a part, moving a cable, pressing a button), billed by the task or the hour, so that the customer need not send its own
staff.
\end{definition}

\section{Cross-connects and the meet-me room}

\begin{definition}[Cross-connect, meet-me room]\label{def:nw:colocation-products-and-how-they-are-sold:mmr}
A \emph{cross-connect}\index{cross-connect} is a cable, usually a pair of fibres, installed and managed by the data-centre operator between
two points of its building: a customer's cabinet and the venue's network, another customer's cabinet, or a telecom carrier's equipment. A
\emph{meet-me room}\index{meet-me room} (MMR) is the area of the building where telecom carriers terminate their networks and from which
cross-connects run to customers.
\end{definition}

Every connection that enters or crosses a \omterm{def:nw:colocation-products-and-how-they-are-sold:colo}{colocation} hall is a \omterm{def:nw:colocation-products-and-how-they-are-sold:mmr}{cross-connect} that someone pays for each month. ICE's operating policies for
Mahwah are typical of the rules: carriers may only sit in the \omterm{def:nw:colocation-products-and-how-they-are-sold:mmr}{meet-me rooms}, with at most 16~kilowatts each; they reach a customer only
through a \omterm{def:nw:colocation-products-and-how-they-are-sold:mmr}{cross-connect}, for which a monthly carrier connection fee is charged; and a customer trading on the venue must use the venue's
network, not a direct \omterm{def:nw:colocation-products-and-how-they-are-sold:mmr}{cross-connect} to the venue's systems. Nasdaq's filings describe the same structure in Carteret: a carrier cage, a
patch panel, cabling managed by the exchange to a distribution point, then to the customer's cabinet, each piece ``color-coded, inventoried,
and auditable''.

\begin{omfigure}
\begin{tikzpicture}[font=\footnotesize,b/.style={draw,rounded corners=2pt,minimum height=0.6cm,align=center,inner sep=3pt},>=Stealth]
\node[b,fill=omThm!15,minimum height=1.8cm] (me) at (0,0) {venue's\\ matching\\ engines and\\ network};
\node[b,fill=omDef!15] (c1) at (3.2,1.9) {near cabinet};
\node[b,fill=omDef!15] (c2) at (6.2,1.9) {middle cabinet};
\node[b,fill=omDef!15] (c3) at (9.2,1.9) {far cabinet};
\draw[thick] (0.95,0.6) -- (3.2,0.6) -- (c1.south);
\draw[thick] (0.95,0.25) -- (6.2,0.25) -- (c2.south);
\draw[thick] (0.95,-0.1) -- (9.2,-0.1) -- (c3.south);
\foreach \r in {0.14,0.19,0.24} {\draw[omMeth,thick,fill=white] (2.2,0.6) circle (\r);}
\foreach \r in {0.14,0.19} {\draw[omMeth,thick,fill=white] (4.6,0.25) circle (\r);}
\node[font=\scriptsize,omMeth] at (2.2,1.05) {long coil};
\node[font=\scriptsize,omMeth] at (4.6,0.7) {short coil};
\node[b,fill=gray!12,minimum height=1.2cm] (mmr) at (0,-2.4) {meet-me room\\ (carriers)};
\node[b,fill=omDef!15] (cage) at (9.2,-1.4) {cage};
\draw[thick,gray] (mmr.north) -- node[left,font=\scriptsize]{carrier links} (me.south);
\draw[thick,gray] (mmr.east) -- (9.2,-2.4) node[pos=0.5,below,font=\scriptsize]{cross-connect from a carrier} -- (cage.south);
\node[font=\scriptsize,align=center] at (6.2,-0.55) {every cable cut to the length of the longest};
\end{tikzpicture}

\omcaption{A venue's data centre, schematically. Customers' cabinets and cages reach the venue's network over cables of equal length; the
nearer a cabinet, the more fibre is coiled in its path. Carriers terminate in the \omterm{def:nw:colocation-products-and-how-they-are-sold:mmr}{meet-me room} and reach customers only through
\omterm{def:nw:colocation-products-and-how-they-are-sold:mmr}{cross-connects} managed by the operator.}\label{fig:nw:colocation-products-and-how-they-are-sold:hall}
\end{omfigure}

\section{Equal cable lengths and the fairness rules}

\begin{definition}[Latency equalisation]\label{def:nw:colocation-products-and-how-they-are-sold:equal}
\emph{Latency equalisation}\index{latency equalisation} is a venue's practice of making every customer's path to its systems equally long
in time, whatever the customer's position in the building, by cutting every connection to the length of the longest (coiling the surplus)
or by adding delay to the shorter ones.
\end{definition}

\begin{proposition}[What a coil costs]\label{prop:nw:colocation-products-and-how-they-are-sold:coil}
If the longest path to the venue's network is $L$ metres of fibre of group index $n_g$, a cabinet whose direct path would be $\ell$ metres
waits $(L - \ell)\,n_g / c_0$ more than it would without equalisation: about \qty{4.88}{\nano\second} per metre for standard fibre.
\end{proposition}

\begin{proof}
Its cable is cut to length $L$; light covers the extra $L-\ell$ metres at $c_0/n_g$.
\end{proof}

A cabinet 5 metres from the engine in a hall whose longest run is 150 metres therefore gives up \qty{707}{\nano\second}, the price of
fairness paid by the customer who would otherwise have won by standing close. The European rule makes the principle law for venues in its
scope, adding that the venue must monitor ``all connections and latency measurements'' and must sell each co-location service on its own
rather than in bundles. In the United States, \omterm{def:nw:colocation-products-and-how-they-are-sold:colo}{colocation} services are exchange services whose fees are filed with the SEC, and exchanges
describe their equalisation in those filings. What equalisation removes is the race for the nearest rack; what it leaves is every other
race: faster servers, faster switches, better code and, for firms that trade across venues, the routes between buildings of chapters~10
to~14.

\section{Port speeds and connectivity products}

A customer's cabinet reaches the venue's systems through the venue's own network, over ports that the venue sells in several speeds and
kinds: market-data ports and order-entry ports, one or ten or more gigabits, a standard network and, at some venues, a lower-latency one
at a higher price, and connections to the venue's disaster-recovery site (chapter~28). The price of a connection is set by filing, like
the cabinet's. The fastest product of \cref{dat:nw:colocation-products-and-how-they-are-sold:fees}, an options exchange's 10-gigabit
ultra-low-latency connection, costs ten times its 1-gigabit connection. Some venues also sell connectivity between data centres
(the wireless links of chapter~14 among them), and the rule of equal conditions applies to what they sell within each product.

\section{Fee schedules and how a footprint is bought}

\begin{dated}{2026-09}{Colocation and connectivity fees, from exchange filings}\label{dat:nw:colocation-products-and-how-they-are-sold:fees}
Nasdaq's filing of October 2024 set, for its NY11-4 expansion hall in Carteret, a monthly fee of 7\,230~USD for an Ultra High Density
Cabinet (above 10 and up to \qty{15}{\kilo\watt}), a cabinet installation fee of 5\,940~USD (cabinet included), and power installation fees of
3\,600~USD (single phase) and 4\,560~USD (three phase); it described its existing monthly cabinet fees as ranging from about 475 to 916~USD
per kilowatt, and its fee increase of November 2024 excluded the NY11-4 fees. Its 2026 filings extend the same installation fees to NY11-5
and list monthly power-circuit fees from 2\,640~USD (single-phase, 20~A, 240~V) to 12\,650.53~USD (three-phase, 32~A, 415~V). MIAX Pearl's
filing of May 2026 raised its monthly fee for a 10-gigabit ultra-low-latency connection to its primary and secondary facilities from
13\,500 to 15\,000~USD, for a 1-gigabit connection from 1\,400 to 1\,500~USD, and for a 10-gigabit connection to its disaster-recovery
facility from 2\,750 to 3\,500~USD.
\end{dated}

\begin{method}[Buying a colocation footprint]\label{met:nw:colocation-products-and-how-they-are-sold:buy}
\begin{enumerate}
\item Size the power first (chapter~8): the servers, switches, timing and capture that must be in the building, at their real draw.
\item Choose cabinets by \omterm{def:nw:colocation-products-and-how-they-are-sold:cabinet}{power density}; a cage when the footprint is large or the firm wants its own access control.
\item List the connections: the venue's market-data and order-entry ports, \omterm{def:nw:colocation-products-and-how-they-are-sold:mmr}{cross-connects} to carriers for the firm's own links, and the
  disaster-recovery connections the venue requires.
\item Price it from the venue's current fee schedule, one-time and monthly, and amortise the one-time fees over the contract's term.
\item Read the venue's operating policies: what may be cross-connected to what, who may enter, remote-hands rates, notice periods.
\item Plan the lead times: cabinets, power and \omterm{def:nw:colocation-products-and-how-they-are-sold:mmr}{cross-connects} take weeks to deliver, and the venue's certification of a new connection
  takes its own.
\end{enumerate}
\end{method}

\begin{omfigure}
\begin{tikzpicture}
\begin{axis}[xbar stacked,width=11.5cm,height=5.8cm,bar width=9pt,xmin=0,xmax=440,
  ytick=data,yticklabels from table={figdata/networks/09-colocation-products-and-how-they-are-sold/bills.csv}{footprint},
  y dir=reverse,xlabel={annual cost (thousand USD)},tick label style={font=\small},label style={font=\small},
  grid=major,grid style={gray!25},legend columns=2,legend style={font=\footnotesize,at={(0.5,-0.22)},anchor=north,draw=none,
  fill=none,/tikz/every even column/.append style={column sep=8pt}},table/col sep=comma]
\addplot[fill=omDef!55,draw=omDef] table[y expr=\coordindex,x=recurring_k]{figdata/networks/09-colocation-products-and-how-they-are-sold/bills.csv};
\addplot[fill=omThm!40,draw=omThm] table[y expr=\coordindex,x=one_time_k]{figdata/networks/09-colocation-products-and-how-they-are-sold/bills.csv};
\legend{monthly fees,one-time fees over 36 months}
\end{axis}
\end{tikzpicture}

\omcaption{Annual cost of six footprints from the published fees of \cref{dat:nw:colocation-products-and-how-they-are-sold:fees}: Nasdaq NY11-4
cabinets (each with installation and three-phase power, the largest with power distribution units) and MIAX Pearl connections. Two
ultra-low-latency connections cost about as much as four cabinets. Data: \texttt{fig\_colo.py} on \texttt{firm.colobill}.}\label{fig:nw:colocation-products-and-how-they-are-sold:bills}
\end{omfigure}

The chart shows where the money goes. A cabinet at Nasdaq costs about 90\,000~USD a year once its installation is spread over three years;
two ultra-low-latency connections to an options exchange cost 360\,000~USD a year, with nothing to install. The building is the cheap
part; the venue's fastest connections are not, and a firm pays for them at every venue it trades.

\omcode{code/firm/colobill/firm_colobill.py}{56}{71}{A footprint's bill: monthly fees, one-time fees amortised over the term, and the date
check that keeps a schedule from going stale.}\label{lst:nw:colocation-products-and-how-they-are-sold:bill}

\section{Tutorial: pricing a footprint}\label{tut:nw:colocation-products-and-how-they-are-sold:tut}

\begin{tutorial}
\textbf{Goal.} Turn two venues' published fees into data, price six footprints and a firm's, and compute what equalisation costs a near
cabinet. \textbf{End state:} \cref{fig:nw:colocation-products-and-how-they-are-sold:bills}, the firm's bill of the problem, and the coil
arithmetic.
\begin{tutsteps}
\item \textbf{Schedules.} \texttt{firm.colobill.SCHEDULES} holds each fee with the ledger row that sources it and the date it was read.
\item \textbf{Footprints.} \texttt{nw\_colo.FOOTPRINTS} lists items and quantities; \texttt{bills()} prices them
  (\cref{lst:nw:colocation-products-and-how-they-are-sold:bill}) and \texttt{python fig\_colo.py} writes the chart's data.
\item \textbf{Staleness.} \texttt{stale(schedule, today)} lists the fees older than 60 days: the dated box's rule, enforced in code.
\item \textbf{The coil.} \texttt{coil(distances, longest)} evaluates \cref{prop:nw:colocation-products-and-how-they-are-sold:coil}.
\end{tutsteps}
\textbf{What to change next.} Add a third schedule from another venue's filing and price the same footprint there; change the term to 60
months and see how little the one-time fees then weigh.
\end{tutorial}

\section{Build: the colocation bill}\label{bld:nw:colocation-products-and-how-they-are-sold:colobill}

\begin{build}
\bfield{Purpose} The firm's \omterm{def:nw:colocation-products-and-how-they-are-sold:colo}{colocation} and connectivity costs computed from published schedules, with provenance and dates, for
chapter~29's annual budget.
\bfield{Interface} \texttt{firm\_colobill}: \texttt{Item(name, monthly, one\_time, unit, source, as\_of)}, \texttt{Schedule(venue, currency,
items)}, \texttt{SCHEDULES}, \texttt{bill(schedule, footprint, term\_months)}, \texttt{stale(schedule, today, days)}, \texttt{break\_even},
\texttt{equalisation\_ns}.
\bfield{Rules} Every fee names its ledger row and date; one-time fees are amortised over a stated term; an unknown item is an error, not a
zero; a fee older than 60 days is reported as stale.
\bfield{Acceptance tests} \texttt{code/firm/colobill/tests/}: a bill's arithmetic by hand, staleness at two dates, break-even, the coil, and an
unknown item refused.
\bfield{Stretch} Read a schedule from a filing's text; add tiered fees (a price that falls with the number of connections).
\end{build}

\begin{omsources}
\item Commission Delegated Regulation (EU) 2017/573 (RTS~10), article~1.
\item SEC Releases 34-101267 (Nasdaq, SR-NASDAQ-2024-056) and 34-105010 (Nasdaq ISE, SR-ISE-2026-09); Federal Register documents 2024-27759,
2026-06357 (Nasdaq) and 2026-10452 (MIAX Pearl).
\item ICE, \emph{ICE Global Network \& Colocation Technical Specifications} and \emph{Mahwah Operating Policies}.
\end{omsources}

\section{Exercises}

\begin{exercise}[$\star$]\label{exo:nw:colocation-products-and-how-they-are-sold:1}
The longest cable in a hall is 120 metres. What delay does equalisation add to a cabinet whose direct path would be 30 metres?
\end{exercise}

\begin{exercise}[$\star$]\label{exo:nw:colocation-products-and-how-they-are-sold:2}
From \cref{dat:nw:colocation-products-and-how-they-are-sold:fees}, what does one NY11-4 Ultra High Density Cabinet cost per kilowatt a month at
10 and at \qty{15}{\kilo\watt}?
\end{exercise}

\begin{exercise}[$\star$]\label{exo:nw:colocation-products-and-how-they-are-sold:3}
What did MIAX Pearl's May 2026 filing do to the annual cost of two 10-gigabit ultra-low-latency connections?
\end{exercise}

\begin{exercise}[$\star\star$]\label{exo:nw:colocation-products-and-how-they-are-sold:4}
Why does a venue sell cabinets by \omterm{def:nw:colocation-products-and-how-they-are-sold:cabinet}{power density} rather than by \omterm{def:nw:servers:ru}{rack units}?
\end{exercise}

\begin{exercise}[$\star\star$]\label{exo:nw:colocation-products-and-how-they-are-sold:5}
A firm proposes to \omterm{def:nw:colocation-products-and-how-they-are-sold:mmr}{cross-connect} its cabinet directly to another firm's cabinet in Mahwah to send it orders. What do ICE's operating
policies say, and why do venues write such rules?
\end{exercise}

\begin{exercise}[$\star\star$]\label{exo:nw:colocation-products-and-how-they-are-sold:6}
Equalisation removes the race for the nearest cabinet. Name three races it leaves.
\end{exercise}

\begin{exercise}[$\star\star\star$]\label{exo:nw:colocation-products-and-how-they-are-sold:7}
\emph{Coding.} With \texttt{firm.colobill}, price four NY11-4 cabinets with installation, three-phase power and PDUs over 36 and over 60
months. By how much does the longer term lower the annual cost?
\end{exercise}

\begin{exercise}[$\star\star\star$]\label{exo:nw:colocation-products-and-how-they-are-sold:8}
\emph{Find the flaw.} ``The venue equalises cable lengths, so every colocated firm sees market data at the same time.''
\end{exercise}

\section{Problem: The Cage That Paid for Itself}

\begin{problem}[{Weekend problem --- a footprint, its bill and its break-even}]\label{pb:nw:colocation-products-and-how-they-are-sold:1}
A market-making firm colocates two Ultra High Density Cabinets in Nasdaq's NY11-4 hall (each with installation and three-phase power) and
buys two 10-gigabit ultra-low-latency connections to an options exchange at MIAX Pearl's fee. It amortises one-time fees over 36 months and
trades 50\,000 contracts a day, 20 days a month.

\medskip\noindent\textbf{Part I --- The bill.}
\begin{enumerate}
\item What are the monthly fees of the two cabinets?
\item What are the one-time fees, and their monthly amortisation?
\item What do the two connections cost a month?
\item What is the footprint's total monthly cost, and its annual cost?
\end{enumerate}

\medskip\noindent\textbf{Part II --- The break-even.}
\begin{enumerate}[resume]
\item How many contracts does the firm trade a month?
\item What must each contract earn, net of every other cost, to pay for the footprint?
\item What if volume halves?
\item Which line of the bill would you cut first if it did?
\end{enumerate}

\medskip\noindent\textbf{Part III --- Fairness and position.}
\begin{enumerate}[resume]
\item The firm's cabinet is 10 metres from the venue's network and the longest run is 150 metres. What does equalisation cost it?
\item Why may the firm not buy a shorter cable?
\item What does the European rule add to that?
\item What does the firm get for its 15\,000~USD a month per connection that a 1-gigabit connection does not give?
\end{enumerate}

\medskip\noindent\textbf{Part IV --- The verdict.}
\begin{enumerate}[resume]
\item State the \emph{named result}: the footprint's monthly all-in cost and the break-even per contract.
\item Over three years, how much of the total is one-time fees?
\item How often must the fee schedule be re-read, and why?
\item What else goes into the footprint that the schedules do not show?
\item How would the bill change at a venue whose connections are cheaper but whose cabinets cost more?
\item Why is the connection, not the cabinet, the expensive part?
\item What would you negotiate, and with whom (One Quant Book~16, chapter~23)?
\item In one sentence: what does a \omterm{def:nw:colocation-products-and-how-they-are-sold:colo}{colocation} footprint buy?
\end{enumerate}
\end{problem}

\section{Interview questions}

\begin{interviewq}[$\star$ \iqroles{developer, trader}]\label{iq:nw:colocation-products-and-how-they-are-sold:1}
What is \omterm{def:nw:colocation-products-and-how-they-are-sold:colo}{colocation}, and what does a firm actually pay for?
\end{interviewq}

\begin{interviewq}[$\star\star$ \iqroles{developer}]\label{iq:nw:colocation-products-and-how-they-are-sold:2}
What is a \omterm{def:nw:colocation-products-and-how-they-are-sold:mmr}{cross-connect} and what is a \omterm{def:nw:colocation-products-and-how-they-are-sold:mmr}{meet-me room}? How does a carrier's circuit reach your cabinet?
\end{interviewq}

\begin{interviewq}[$\star\star$ \iqroles{trader, developer}]\label{iq:nw:colocation-products-and-how-they-are-sold:3}
Why do venues equalise cable lengths, and what advantage remains for a colocated firm?
\end{interviewq}

\begin{interviewq}[$\star\star$ \iqroles{developer}]\label{iq:nw:colocation-products-and-how-they-are-sold:4}
You are moving into a venue's data centre in three months. What do you order, in what order?
\end{interviewq}

\begin{interviewq}[$\star\star\star$ \iqroles{trader}]\label{iq:nw:colocation-products-and-how-they-are-sold:5}
How would you decide whether a \omterm{def:nw:colocation-products-and-how-they-are-sold:colo}{colocation} footprint at a venue pays for itself?
\end{interviewq}

\begin{interviewq}[$\star\star$ \iqroles{developer}]\label{iq:nw:colocation-products-and-how-they-are-sold:6}
Where do you find a U.S. exchange's \omterm{def:nw:colocation-products-and-how-they-are-sold:colo}{colocation} fees, and how do you know they are current?
\end{interviewq}
